% META: {"id":"1SPE-DERLOCAL-EX-045","chapitre":"1SPE-DERIVATION-LOCAL","type_objet":"exercice","capacites":["1SPE-DERIVATION-LOCAL-C5"],"capacites_codes":["C5"],"methodes":["M1"],"parcours":2,"competences":["calculer","modeliser"],"duree_min":12,"mode_creation":"ex_nihilo","sources_inspiration":[],"fichier_tex":"chapitres/1SPE-DERIVATION-LOCAL/exercices/1SPE-DERLOCAL-EX-045.tex","corrige_tex":"chapitres/1SPE-DERIVATION-LOCAL/corriges/1SPE-DERLOCAL-CO-045.tex","status":"generated"}
% BEGIN-VERIFY
% from sympy import symbols, Rational, Abs
% x = symbols('x')
% f = (1 + x)**5
% assert f.subs(x, 0) == 1
% fp = f.diff(x)
% assert fp.subs(x, 0) == 5
% approx = 1 + 5*Rational(2, 100)
% assert approx == Rational(11, 10)
% exact = Rational(102, 100)**5
% assert exact == Rational(11040808032, 10000000000)
% END-VERIFY
\begin{exercice}{1SPE-DERLOCAL-EX-045}{2}{12}
On souhaite obtenir une valeur approchée de $(1{,}02)^5$ sans calculatrice, en utilisant l'approximation affine.

On pose $f(x) = (1+x)^5$.
\begin{enumerate}
\item Calculer $f(0)$ et $f'(0)$. On admet que $f'(x) = 5(1+x)^4$.
\item Écrire l'approximation affine de $f$ au voisinage de $x = 0$ : $f(x) \approx f(0) + f'(0) \cdot x$.
\item En déduire une valeur approchée de $(1{,}02)^5$ en prenant $x = 0{,}02$.
\item Comparer avec la valeur exacte $(1{,}02)^5 = 1{,}104\,080\,8\ldots$ Commenter la précision.
\end{enumerate}
\end{exercice}
